%----------------------------------------------------------------------------------------
% METADATA BLOCK (YAML)  — see inbox/README.md
%----------------------------------------------------------------------------------------
% uid: ""
% school: "lycee"
% level: "3eme"
% section: "Economie et Gestion"
% chapter: "Statistiques"
% title: "Réclamations journalières"
% difficulty: "Facile"
% tags: ["statistiques", "moyenne", "variance", "ecart-type"]
% provenance: "own"
% batch: "2026-11-19-3eco-statistiques-n5"
%----------------------------------------------------------------------------------------

\newcommand{\R}{\mathbb{R}}
\newcommand{\vect}[1]{\overrightarrow{#1}}

% ---STATEMENT---
Un service après-vente a noté le nombre de réclamations reçues chaque jour pendant $20$ jours :
\begin{center}
\begin{tabular}{|l|c|c|c|c|c|}\hline
Réclamations $x_i$ & $1$ & $2$ & $3$ & $4$ & $5$ \\\hline
Nombre de jours $n_i$ & $2$ & $5$ & $6$ & $5$ & $2$ \\\hline
\end{tabular}\end{center}
\begin{enumerate}
  \item Calculer le nombre moyen de réclamations par jour.
  \item Calculer la variance $V$ et l'écart type $\sigma$ (au centième).
  \item Quel pourcentage des jours compte au moins $4$ réclamations ?
\end{enumerate}

% ---SOLUTION---
\begin{enumerate}[itemsep=14pt]
  \item $N=20$ :
  \[\overline x=\dfrac{2\times1+5\times2+6\times3+5\times4+2\times5}{20}=\dfrac{2+10+18+20+10}{20}=\dfrac{60}{20}=3\]

  \item
  \begin{prop}[Variance et écart type]
  \ib{$V=\dfrac{n_1x_1^2+\dots+n_px_p^2}{N}-\overline x^{\,2}$} \quad et \quad \ib{$\sigma=\sqrt V$}.
  \end{prop}
  On calcule d'abord la moyenne des carrés :
  \[\dfrac{2\times1+5\times4+6\times9+5\times16+2\times25}{20}=\dfrac{2+20+54+80+50}{20}=\dfrac{206}{20}=10{,}3\]
  \[V=10{,}3-3^2=10{,}3-9=1{,}3\qquad\sigma=\sqrt{1{,}3}\approx1{,}14\]

  \item Au moins $4$ réclamations : $5+2=7$ jours sur $20$, soit $\dfrac{7}{20}=35\,\%$.
\end{enumerate}
